\chapter{Types}
\label{ch:types}

\haira{} is a statically-typed language with structural typing and type inference for local variables. All types are known at compile time.

\section{Type System Overview}

\begin{itemize}
    \item \textbf{Static typing} -- All types resolved at compile time
    \item \textbf{Structural typing} -- Type compatibility based on structure, not name
    \item \textbf{Type inference} -- Local variables infer types from initializers
    \item \textbf{Explicit annotations} -- Required for function signatures
    \item \textbf{No null} -- Option types (\type{T?}) replace null pointers
\end{itemize}

\section{Primitive Types}

\subsection{Integer Types}

\begin{table}[h]
\centering
\begin{tabular}{lll}
\toprule
\textbf{Type} & \textbf{Size} & \textbf{Range} \\
\midrule
\type{int} & 64 bits & $-2^{63}$ to $2^{63}-1$ \\
\type{i8} & 8 bits & $-128$ to $127$ \\
\type{i16} & 16 bits & $-32768$ to $32767$ \\
\type{i32} & 32 bits & $-2^{31}$ to $2^{31}-1$ \\
\type{i64} & 64 bits & $-2^{63}$ to $2^{63}-1$ \\
\type{u8} & 8 bits & $0$ to $255$ \\
\type{u16} & 16 bits & $0$ to $65535$ \\
\type{u32} & 32 bits & $0$ to $2^{32}-1$ \\
\type{u64} & 64 bits & $0$ to $2^{64}-1$ \\
\bottomrule
\end{tabular}
\caption{Integer types}
\end{table}

\begin{lstlisting}
count: int = 42
byte: u8 = 255
offset: i32 = -100
\end{lstlisting}

\begin{note}
The default \type{int} is an alias for \type{i64}. Integer literals without explicit type default to \type{int}.
\end{note}

\subsection{Floating-Point Types}

\begin{table}[h]
\centering
\begin{tabular}{lll}
\toprule
\textbf{Type} & \textbf{Size} & \textbf{Precision} \\
\midrule
\type{float} & 64 bits & IEEE 754 double \\
\type{f32} & 32 bits & IEEE 754 single \\
\type{f64} & 64 bits & IEEE 754 double \\
\bottomrule
\end{tabular}
\caption{Floating-point types}
\end{table}

\begin{lstlisting}
pi: float = 3.14159
temperature: f32 = 98.6
\end{lstlisting}

\subsection{Boolean Type}

\begin{lstlisting}
is_valid: bool = true
is_empty: bool = false
\end{lstlisting}

Boolean values are used in conditionals and logical expressions. Only \code{true} and \code{false} are valid boolean values---there is no implicit conversion from other types.

\subsection{String Type}

Strings are immutable sequences of UTF-8 encoded bytes:

\begin{lstlisting}
name: string = "Haira"
greeting: string = "Hello, " + name
\end{lstlisting}

String operations:
\begin{lstlisting}
import "string"

s = "Hello, World!"
length = string.len(s)        // 13
upper = string.to_upper(s)    // "HELLO, WORLD!"
parts = string.split(s, ", ") // ["Hello", "World!"]
\end{lstlisting}

\section{Composite Types}

\subsection{Arrays}

Arrays are ordered, homogeneous, dynamically-sized collections:

\begin{lstlisting}
// Array literal
numbers: []int = [1, 2, 3, 4, 5]

// Empty array
empty: []string = []

// Array operations
import "array"

first = numbers[0]            // 1
length = array.len(numbers)   // 5
numbers = array.push(numbers, 6)
\end{lstlisting}

\begin{lstlisting}[language=EBNF]
array_type = "[" "]" type ;
\end{lstlisting}

\subsection{Maps}

Maps are unordered key-value collections:

\begin{lstlisting}
// Map literal
scores: [string:int] = [
    "alice": 100,
    "bob": 95,
    "charlie": 87
]

// Empty map
empty: [string:User] = [:]

// Map operations
alice_score = scores["alice"]    // 100
scores["david"] = 92             // Insert
\end{lstlisting}

\begin{lstlisting}[language=EBNF]
map_type = "[" type ":" type "]" ;
\end{lstlisting}

\begin{note}
Map keys must be comparable types: primitives, strings, or structs containing only comparable fields.
\end{note}

\subsection{Tuples}

Tuples are fixed-size, heterogeneous collections:

\begin{lstlisting}
// Tuple type
point: (int, int) = (10, 20)

// Accessing elements
x = point.0  // 10
y = point.1  // 20

// Multiple return values use tuples
fn divide(a: int, b: int) -> (int, int) {
    return a / b, a % b
}

quotient, remainder = divide(17, 5)
\end{lstlisting}

\begin{lstlisting}[language=EBNF]
tuple_type = "(" type { "," type } ")" ;
\end{lstlisting}

\section{Option Types}

Option types represent values that may or may not exist, replacing null pointers:

\begin{lstlisting}
// Optional value
maybe_name: string? = nil
maybe_age: int? = 25

// Checking for nil
if maybe_name != nil {
    io.println(maybe_name)
}

// Default value with 'or'
name = maybe_name or "Unknown"
\end{lstlisting}

\begin{lstlisting}[language=EBNF]
option_type = type "?" ;
\end{lstlisting}

\subsection{Unwrapping Options}

\begin{lstlisting}
value: int? = get_value()

// Safe unwrap with default
result = value or 0

// Conditional unwrap
if value != nil {
    // value is automatically unwrapped here
    io.println(value)
}

// Match on option
match value {
    nil => io.println("No value")
    v => io.println("Got: " + to_string(v))
}
\end{lstlisting}

\section{Struct Types}

Structs are named product types with named fields:

\begin{lstlisting}
struct User {
    id: int
    name: string
    email: string
    active: bool
}

// Creating a struct
user = User{
    id: 1,
    name: "Alice",
    email: "alice@example.com",
    active: true
}

// Accessing fields
io.println(user.name)

// Updating fields (structs are mutable by default)
user.active = false
\end{lstlisting}

\subsection{Struct Methods}

Methods can be defined on structs:

\begin{lstlisting}
struct Rectangle {
    width: float
    height: float
}

fn (r: Rectangle) area() -> float {
    return r.width * r.height
}

fn (r: Rectangle) perimeter() -> float {
    return 2 * (r.width + r.height)
}

// Usage
rect = Rectangle{width: 10.0, height: 5.0}
io.println(rect.area())       // 50.0
io.println(rect.perimeter())  // 30.0
\end{lstlisting}

\subsection{Nested Structs}

\begin{lstlisting}
struct Address {
    street: string
    city: string
    country: string
}

struct Person {
    name: string
    address: Address
}

person = Person{
    name: "Bob",
    address: Address{
        street: "123 Main St",
        city: "New York",
        country: "USA"
    }
}

io.println(person.address.city)  // "New York"
\end{lstlisting}

\section{Enum Types}

Enums define a type with a fixed set of named values:

\begin{lstlisting}
enum Color {
    Red,
    Green,
    Blue
}

enum Status {
    Pending,
    Active,
    Completed,
    Failed
}

// Usage
color: Color = Color.Red
status: Status = Status.Active

match color {
    Color.Red => io.println("Stop")
    Color.Green => io.println("Go")
    Color.Blue => io.println("Info")
}
\end{lstlisting}

\subsection{Enums with Associated Data}

Enums can carry associated data (tagged unions):

\begin{lstlisting}
enum Result<T> {
    Ok(T),
    Err(Error)
}

enum Message {
    Text(string),
    Number(int),
    Pair(int, int)
}

// Usage
msg: Message = Message.Text("Hello")

match msg {
    Message.Text(s) => io.println("Text: " + s)
    Message.Number(n) => io.println("Number: " + to_string(n))
    Message.Pair(a, b) => io.println("Pair: " + to_string(a) + ", " + to_string(b))
}
\end{lstlisting}

\section{Type Aliases}

Type aliases create alternative names for existing types:

\begin{lstlisting}
type UserId = int
type UserMap = [string:User]
type Handler = fn(Request) -> Response

// Usage
id: UserId = 42
users: UserMap = [:]
\end{lstlisting}

\section{Function Types}

Functions are first-class values with types:

\begin{lstlisting}
// Function type syntax
type Comparator = fn(int, int) -> bool
type Transformer = fn(string) -> string
type Callback = fn()

// Higher-order functions
fn apply(f: fn(int) -> int, x: int) -> int {
    return f(x)
}

fn double(n: int) -> int {
    return n * 2
}

result = apply(double, 21)  // 42
\end{lstlisting}

\section{Generic Types}

Generic types parameterize over other types:

\begin{lstlisting}
// Generic struct
struct Pair<T, U> {
    first: T
    second: U
}

// Usage
pair: Pair<int, string> = Pair{first: 1, second: "one"}

// Generic function
fn swap<T, U>(p: Pair<T, U>) -> Pair<U, T> {
    return Pair{first: p.second, second: p.first}
}

swapped = swap(pair)  // Pair<string, int>
\end{lstlisting}

\subsection{Type Constraints}

Generics can be constrained to types implementing certain behaviors:

\begin{lstlisting}
// Constraint definition
constraint Comparable {
    fn compare(other: Self) -> int
}

// Constrained generic
fn max<T: Comparable>(a: T, b: T) -> T {
    if a.compare(b) > 0 {
        return a
    }
    return b
}

// Multiple constraints
fn process<T: Comparable + Printable>(value: T) {
    io.println(value.to_string())
}
\end{lstlisting}

\section{The Error Type}

The built-in \type{Error} type represents errors:

\begin{lstlisting}
struct Error {
    message: string
    code: int?
    source: Error?
}

// Creating errors
err = Error{message: "file not found"}
err_with_code = Error{message: "permission denied", code: 403}

// Error chaining
wrapped = Error{
    message: "failed to read config",
    source: err
}
\end{lstlisting}

\section{Type Inference}

\haira{} infers types for local variables from their initializers:

\begin{lstlisting}
// Types inferred
x = 42              // int
pi = 3.14           // float
name = "Haira"      // string
active = true       // bool
numbers = [1, 2, 3] // []int

// Explicit annotation when needed
count: i32 = 100    // Specific integer type
\end{lstlisting}

\begin{warning}
Type annotations are \textbf{required} for:
\begin{itemize}
    \item Function parameters
    \item Function return types
    \item Struct fields
    \item Uninitialized variables
\end{itemize}
\end{warning}

\section{Structural Typing}

Types are compatible based on structure, not name:

\begin{lstlisting}
struct Point2D {
    x: float
    y: float
}

struct Vector2D {
    x: float
    y: float
}

fn magnitude(p: {x: float, y: float}) -> float {
    return math.sqrt(p.x * p.x + p.y * p.y)
}

// Both work because they have the same structure
point = Point2D{x: 3.0, y: 4.0}
vector = Vector2D{x: 3.0, y: 4.0}

io.println(magnitude(point))   // 5.0
io.println(magnitude(vector))  // 5.0
\end{lstlisting}

\section{Type Grammar}

\begin{lstlisting}[language=EBNF]
type = primitive_type
     | array_type
     | map_type
     | tuple_type
     | option_type
     | function_type
     | generic_type
     | struct_type
     | identifier ;

primitive_type = "int" | "i8" | "i16" | "i32" | "i64"
               | "u8" | "u16" | "u32" | "u64"
               | "float" | "f32" | "f64"
               | "bool" | "string" ;

array_type = "[" "]" type ;

map_type = "[" type ":" type "]" ;

tuple_type = "(" type { "," type } ")" ;

option_type = type "?" ;

function_type = "fn" "(" [ type_list ] ")" [ "->" type ] ;

generic_type = identifier "<" type_list ">" ;

type_list = type { "," type } ;
\end{lstlisting}
